La Parker Solar Probe és una nau espacial destinada a acostar-se molt a la superfície solar. La gràfica següent ens mostra com varia la distància al Sol al llarg dels primers 1000 dies de missió i s'hi indiquen els instants A, B i C. Les unitats de la distància a la superfície del Sol són en radis solars, $R_\mathrm{S}$.

\figura[0.85]{fig1.png}

\begin{apartat}{1}
Observeu a la gràfica els moments de màxim acostament al Sol de cada òrbita i determineu quantes voltes complertes ha donat al voltant del Sol en aquests 1000 dies. Quina és la mida de l'eix major de l'òrbita entre els moments A i C (doneu el resultat en radis solars)?
\end{apartat}

\begin{apartat}{1}
Representeu esquemàticament el Sol i l'òrbita de la nau entre els moments A i C. Indiqueu sobre el dibuix les posicions corresponents a A, B i C. Situa la nau a la posició B i dibuixa-hi en aquest instant els vectors velocitat i acceleració de la nau (no cal calcular-ne els mòduls). En quina posició la velocitat de la nau és màxima? Justifiqueu la resposta indicant el principi físic en que us baseu. En quina posició la velocitat de la nau és màxima? Justifiqueu la resposta utilitzant el principi de conservació del moment angular.
\end{apartat}
